\begin{claim}{D.3.6}{Proved}{The Moore target gives the sphere test in the primary interval}
Fix a prime $p$ and $d\geq3$. For every pointed source $X$,
\[
 \Map_*(X,P^d(p))\simeq *
 \quad\Longrightarrow\quad
 \Map_*(X,(S^n)^{\wedge}_p)\simeq *
                       \quad(n\geq2).
 \tag{D.3.6.1}
\]
The reverse implication holds if
$\widetilde H_*(X;\QQ)=0$. In particular, if
$\widetilde H_*(X;\ZZ[1/p])=0$, the single Moore target and
the whole family of simply connected $p$-completed spheres
have the same nullity test. For $p=2$, the localizer furnished
by the Moore wedges of D.3.4 is therefore exactly
\[
                 W_S\cap R(\loc(\mathbb S/2)).
 \tag{D.3.6.2}
\]
This identification does not establish that it is strictly
larger than $W_{\mathrm{rev}}$.
\end{claim}
\begin{proof}
Fix $X$ satisfying the left side of (D.3.6.1) and put
$\mathcal R_X=\{Z:\Map_*(X,Z)\simeq *\}$.
We give both the higher-exponent construction and the
unstable inverse-limit justification explicitly.

\medskip\noindent
\textbf{A smash-closed resolving family.}
Put $M=P^3(p)$ and
\[
 \mathcal E=\{\Sigma^i M^{\wedge q}:i\geq0,\ q\geq1\}.
\]
This family is closed under suspension and smash product,
and its members are simply connected. We claim that every
member of $\mathcal E$ belongs to $\mathcal R_X$.

The constructions of D.1.7--D.1.8 supply arbitrarily large
integers $D$ for which
$\Omega(P^D(p)\vee P^D(p))\in\mathcal R_X$.
For odd $p$, use the finite-bouquet conclusion of D.1.7.
For $p=2$, use the iterated two-copy retraction of D.1.8;
its upper dimensions increase without bound.
Set $K=P^{D-1}(p)=\Sigma^{D-4}M$. The two-letter
Hilton--Milnor argument then gives
$\Omega\Sigma K^{\wedge m}\in\mathcal R_X$ for all $m\geq1$.
For fixed $q$, apply Hilton--Milnor to
$\Sigma K^{\wedge q}\vee\Sigma K^{\wedge q}$.
Its factors have the form $\Omega\Sigma K^{\wedge qj}$,
and their connectivities tend to infinity, so its loop space
also belongs to $\mathcal R_X$. The two-copy
James-resolution argument of D.1.8, with $K^{\wedge q}$ as
one letter, now puts every finite bouquet of $K^{\wedge q}$
in $\mathcal R_X$.

Since
$K^{\wedge q}=\Sigma^{q(D-4)}M^{\wedge q}$,
the finite cone cubes of \cite[Theorem~8]{I_Strom2012}
descend this entire finite-bouquet family to
$\Sigma^iM^{\wedge q}$ whenever $0\leq i\leq q(D-4)$.
There are arbitrarily large choices of $D$, so the asserted
membership holds for every $i,q$.
Theorem~7 of the same paper now gives
$\Sigma\mathcal E^\vee\subseteq\mathcal R_X$, where the
superscript denotes finite-type wedges. Because
$\mathcal R_X$ is a strong resolving class, its Corollary~11
puts in $\mathcal R_X$ every simply connected space with
finite $\mathcal E^\vee$-cone length. No claim of automatic
closure under arbitrary cofiber sequences is being made.

\medskip\noindent
\textbf{The exponent can be increased.}
For $n\geq3$ and $r\geq1$, consider the actual composite
\[
 P^n(p)\xrightarrow{q}S^n
       \xrightarrow{\iota}P^{n+1}(p^r),
 \tag{D.3.6.3}
\]
where $q$ pinches the bottom cell and $\iota$ includes the
bottom sphere. Its cofiber has two cells in degree $n$ and
two in degree $n+1$, with integral cellular boundary, after
choosing orientations,
\[
                    \begin{pmatrix}p^r&1\\0&p\end{pmatrix}.
\]
Changing the bases of the two wedges of $n$-spheres and
cancelling the degree-one cell pair leaves the attaching
degree $p^{r+1}$. This is an unstable cellular argument:
for $n\geq2$, maps from an $n$-sphere to a wedge of
$n$-spheres are classified by their degrees. Hence the
cofiber of (D.3.6.3) is $P^{n+1}(p^{r+1})$.

Starting from the cofiber of $P^n(p)\to *$, this constructs
$P^{n+1}(p^r)$ with finite $\mathcal E$-cone length.
The same holds for every finite bouquet of these spaces.
The preceding resolving-class result makes all such bouquets
null. The Barratt cone cubes lower the bouquets with top
dimension four to top dimension three, if necessary.
Consequently
\[
               P^{n+1}(p^r)\in\mathcal R_X
                         \quad(n\geq2,\ r\geq1).
 \tag{D.3.6.4}
\]

\medskip\noindent
\textbf{The unstable Moore tower.}
For $n\geq2$, put $M_r=P^{n+1}(p^r)$ and use the transition
$M_{r+1}\to M_r$ induced by the commuting square of sphere
attaching maps whose top degree is $p$ and whose bottom
degree is $1$. The bottom-sphere inclusions are compatible.
We prove
\[
                   \holim_r M_r\simeq(S^n)^{\wedge}_p.
 \tag{D.3.6.5}
\]
The assertion concerns spaces, not the inverse limit of
Moore spectra.

Here are the precise pro-space tools. In
\cite[Theorems~3.3, 6.3, 6.7 and~7.2; Proposition~7.3]
{I_Isaksen2005}, a pro-space is fibrant for the
$\FF_p$-cohomology model structure if it is strictly fibrant
and essentially levelwise
$\{K(\FF_p,j):j\geq0\}$-nilpotent. Proposition~3.5 there
identifies this last condition with a finite tower of
principal fibrations with those fibres. The
$\FF_p$-homology structure has fewer acyclic cofibrations,
so an object fibrant for the cohomology structure is
also fibrant for the homology structure. Proposition~7.3
identifies a homological fibrant replacement of a connected
space with the Postnikov-truncated Bousfield--Kan tower.
The cited numbering is that of the opened author preprint.

Apply these facts to the pro-space
\[
                    Y_{r,N}=P_NM_r\qquad(r,N\geq1).
\]
Each of its levels is simply connected and has only
finitely many nonzero homotopy groups, all finite
$p$-groups. Its Postnikov tower is a finite principal
tower with fibres $K(A,j)$ for finite abelian $p$-groups
$A$. Refining each $A$ by a finite filtration with
successive quotient $\FF_p$ refines this to a finite
principal tower with fibres $K(\FF_p,j)$. Thus every level
is nilpotent in the precise sense required above.
A strict fibrant replacement $Y'$ remains essentially
levelwise of this type, and so is fibrant in both of
Isaksen's model structures. Equivalently, choose a
cofinal diagonal tower and its ordinary fibration
replacement; each replacement level is weakly equivalent
to the original finite Postnikov level.

The compatible map $cS^n\to Y$ is an isomorphism on
pro-homology with coefficients in $\FF_p$. Indeed for
each fixed degree, sufficiently high Postnikov truncation
does not change that homology. The Moore space has one
reduced mod-$p$ class in degree $n$ and one in degree
$n+1$. The first has identity transition, while the second
has zero transition because the top cellular degree is $p$.
The second pro-group is therefore zero. This verifies the
homological weak-equivalence hypothesis, including
coefficients in every $\FF_p$-module.

It follows that $Y'$ is a homological fibrant model of the
completion of $S^n$. By Proposition~7.3 and Theorem~7.2
of the cited source, it is weakly equivalent between
fibrant objects to a strict fibrant model of
$\{P_NR_NS^n\}_N$, where $R_NS^n$ is the $N$th
Bousfield--Kan $\FF_p$-tower stage.
Taking the derived inverse limit preserves that
equivalence. To check this model-categorically, the
constant-pro-space functor preserves cofibrations and
ordinary weak equivalences, and hence is left Quillen
from spaces to the homological pro-space structure;
its right adjoint is inverse limit.
The derived limit of the latter tower is
$\holim_NR_NS^n=(S^n)^{\wedge}_p$: removing the
Postnikov truncations does not change the limit, since
in each fixed homotopy degree the truncation is eventually
an isomorphism, also one degree above it for the
$\lim^1$ term.
On the other hand,
\[
 \holim_{r,N}P_NM_r
       \simeq\holim_r\bigl(\holim_NP_NM_r\bigr)
       \simeq\holim_rM_r.
 \]
This proves (D.3.6.5). Combining it with (D.3.6.4) and
closure under target homotopy limits proves (D.3.6.1)
for arbitrary sources.

\medskip\noindent
\textbf{The reverse test on rationally acyclic sources.}
Suppose now that $X$ is rationally acyclic and null for
all $(S^n)^{\wedge}_p$, $n\geq2$. The arithmetic square
for each finite sphere identifies its mapping space
from $X$ to $S^n_{(p)}$ with that to $(S^n)^{\wedge}_p$:
the other two vertices are rational and have contractible
mapping spaces from $X$. This is the same direct
arithmetic-square argument as in C.8 and G.6, with $p$
in place of $2$. Thus $S^n_{(p)}\in\mathcal R_X$ for
all $n\geq2$.

The family of simply connected $p$-local spheres is
closed under smash products and suspension. Theorems~7
and~10 and Corollary~11 of \cite{I_Strom2012} therefore
put in $\mathcal R_X$ every sufficiently suspended
$p$-local finite simply connected complex, and every
finite bouquet of its sufficiently suspended copies:
its finite skeletal cone decomposition has attaching
spaces that are finite wedges of such spheres.
Localization preserves these cofiber, suspension and
wedge comparisons by
\cite[Propositions~6.4.1--6.4.6]{FGHMayPonto2012}.
In particular take finite bouquets of $\Sigma^2P^d(p)$.
The finite cone cubes descend them twice and give
$P^d(p)\in\mathcal R_X$.
This proves the reverse implication under the stated
rational hypothesis.

Finally every cofiber in $R(\loc(\mathbb S/2))$ is
rationally acyclic. On such cofibers G.6 identifies the
completed-sphere test with the localized-sphere test
defining $W_S$. The cofiber characterizations of the
localizers now give (D.3.6.2).
\end{proof}
